\section{Introduction}

\subsection{Objet du document}

Le présent document constitue la première analyse fonctionnelle de l'application de suivi des dépenses budgétaires. Il formalise et structure les éléments recueillis lors de la réunion initiale organisée avec le responsable du projet.

Son objectif est de fournir à la maîtrise d'ouvrage une représentation claire, partagée et vérifiable du fonctionnement attendu du futur système, avant d'engager les travaux de modélisation des données et de conception technique.

\subsection{Démarche méthodologique et périmètre d'analyse}

Conformément aux principes d'ingénierie des systèmes d'information administratifs et financiers, cette première étape se concentre exclusivement sur la vision fonctionnelle externe du système :

\begin{itemize}
  \item la structuration du processus métier de bout en bout, depuis l'expression du besoin jusqu'à la transmission du dossier d'ordonnancement au comptable public ;
  \item l'identification précise des acteurs intervenant dans la chaîne de la dépense et la délimitation de leurs responsabilités respectives ;
  \item la définition des fonctionnalités attendues du système à travers un diagramme général des cas d'utilisation conforme à la norme UML ;
  \item le recensement des zones d'incertitude et des règles de gestion qui devront être arbitrées par le responsable du projet.
\end{itemize}

\begin{encadre}{Principe de rigueur documentaire}
Afin de garantir la fidélité de l'analyse aux besoins exprimés, ce rapport distingue strictement :
\begin{enumerate}
  \item \textbf{les faits établis}, directement issus des notes de la réunion de cadrage ;
  \item \textbf{les hypothèses de structuration}, introduites uniquement pour organiser logiquement l'enchaînement des opérations ;
  \item \textbf{les points à valider}, qui font l'objet d'un questionnaire dédié en section 8 et ne doivent pas être considérés comme des règles métier définitives.
\end{enumerate}
Aucune responsabilité ni étape supplémentaire n'a été extrapolée au-delà des éléments communiqués.
\end{encadre}
